\documentclass[11pt]{article}
\usepackage[margin=1in]{geometry}
\usepackage{amsmath,amssymb,amsthm}
\usepackage{microtype}
\usepackage{hyperref}
\hypersetup{colorlinks=true,linkcolor=black,citecolor=black,urlcolor=blue,
  pdftitle={Irreducibility of exposed-point varieties of discotopes},
  pdfauthor={Alec Kriebel}}
\newtheorem{theorem}{Theorem}
\newtheorem{lemma}[theorem]{Lemma}
\newcommand{\Exp}{\operatorname{Exp}}
\newcommand{\Span}{\operatorname{span}}
\newcommand{\R}{\mathbb R}
\newcommand{\C}{\mathbb C}
\title{Irreducibility of exposed-point varieties of discotopes}
\author{Alec Kriebel\thanks{Independent researcher. ORCID:
  \href{https://orcid.org/0009-0001-9320-500X}{0009-0001-9320-500X}.}}
\date{September 30, 2026\quad (version 1.0)}
\begin{document}
\maketitle
\begin{abstract}
We prove that the complex Zariski closure of the exposed points of any
finite Minkowski sum of ellipsoidal discs of dimension at least two is
irreducible. No genericity or full-dimensionality assumption is needed.
This proves the exposed-point irreducibility conjecture stated by Meroni
in the 2023 Oberwolfach report. Under an explicit generic condition on
the summand spans, we also prove that this closure equals the purely
nonlinear part defined by Gesmundo and Meroni, thereby proving their
2022 irreducibility conjecture. The proof uses a connected real-analytic
support-point parametrization. A nongeneric example shows why the two
varieties must be distinguished.
\end{abstract}
\noindent\textbf{Status.} Unrefereed preprint. Generative AI was used
extensively in developing and checking the result and preparing this
manuscript; no human peer review or formal proof certification is claimed.

\section{The conjectures and the result}

A \emph{disc} is a linear image of a closed Euclidean unit ball. Thus
it is an ellipsoidal ball in its linear span and is centered at the
origin. A \emph{discotope} is a finite Minkowski sum
\(D=D_1+\cdots+D_N\subset\R^d\), with \(N\ge1\).
A point of a compact convex set \(K\) is \emph{exposed} if it uniquely
maximizes some nonzero real linear functional on \(K\). Write
\[
 E(K)=\overline{\Exp(K)}^{\mathrm{Zar},\C}\subset\C^d.
\]
Meroni's Conjecture~1 in \cite[p.~830]{Meroni2023} asks whether \(E(D)\)
is irreducible for generic discotopes whose summands have dimension at
least two. We prove the following stronger statement.

\begin{theorem}\label{thm:main}
Let \(D=D_1+\cdots+D_N\subset\R^d\) be a discotope, with every
\(D_i\) of dimension at least two. Then \(E(D)\) is irreducible over
\(\C\), without genericity or full-dimensionality assumptions.
\end{theorem}

Gesmundo and Meroni proved irreducibility for generic discotopes in
the range \(\sum_i(\dim D_i-1)\le d-1\) (Theorem~4.3), and for
generic sums of two-dimensional discs (Theorems~4.3 and~6.1).
Remark~3.4 treats sums whose summands are all full-dimensional.
Their Section~5 already uses the support-gradient parametrization.
The contribution here is the general irreducibility conclusion for
arbitrary summand ranks at least two, together with the generic
identification of the two varieties. The proof uses elementary
analytic and convex facts rather than an asserted irreducibility
theorem for a complex critical locus.

There is a distinction between this target and the earlier formulation.
Gesmundo and Meroni define the \emph{purely nonlinear part}
\cite[Definition~3.3]{GM2022} by
\begin{equation}\label{eq:S}
 S(D)=\overline{D^\partial\cap\partial D}^{\mathrm{Zar},\C},
 \qquad D^\partial=\sum_i\partial_{L_i}D_i,
 \qquad L_i=\Span D_i,
\end{equation}
where each summand boundary is relative to its own span. Their
Conjecture~8.2 asks for \(S(D)\) to be irreducible for generic
discotopes with no one-dimensional summands. In Section~\ref{sec:generic}
we prove \(S(D)=E(D)\) under an explicit condition holding generically,
which proves that conjecture too. The two varieties can differ outside
generic position, as Section~\ref{sec:limits} demonstrates. We do not
address the separate questions about degrees, birationality, or the
entire complex critical locus of the addition map.

\section{A connected analytic parametrization}

We first record two elementary facts. The analytic fact is a standard
identity-theorem argument. Kucharz and Kurdyka explicitly record its
real Zariski version in the proof of Proposition~4.3 of
\cite[p.~95]{KK2017}. We include the complex version for completeness.

\begin{lemma}\label{lem:connected}
The complement in \(\R^d\) of a finite union of real linear subspaces
of codimension at least two is nonempty, open, and path connected.
\end{lemma}
\begin{proof}
A finite union of proper real subspaces cannot fill \(\R^d\): choose a
nonzero linear form vanishing on each subspace; their product is a
nonzero polynomial. The complement is open because the union is closed.
Let \(x,y\) be in the complement. If the excluded subspaces are \(M_i\),
each \(M_i+\R x\) and \(M_i+\R y\) has dimension at most \(d-1\).
Choose \(z\) outside their finite union. If
\((1-t)x+tz\in M_i\) for \(0<t\le1\), then
\(z\in M_i+\R x\), a contradiction. The same holds for the segment
from \(y\) to \(z\). These two segments give the required path.
\end{proof}

\begin{lemma}\label{lem:analytic}
If \(U\subset\R^d\) is nonempty, open, and connected and
\(F:U\to\R^m\) is real analytic, the complex Zariski closure of
\(F(U)\) is irreducible.
\end{lemma}
\begin{proof}
Let \(I\subset\C[X_1,\ldots,X_m]\) be the ideal of polynomials
vanishing on \(F(U)\). It is proper. If \(PQ\in I\) and \(P\notin I\),
there is a nonempty open subset of \(U\) on which \(P\circ F\ne0\),
so \(Q\circ F=0\) there. The real-analytic identity theorem, applied
to the real and imaginary parts on connected \(U\), implies
\(Q\circ F=0\) throughout \(U\). Hence \(I\) is prime. It is exactly
the vanishing ideal of the complex Zariski closure, so that closure is
irreducible.

For completeness, analytic uniqueness here requires only connected
openness. The points where a real-analytic function vanishes on a
neighborhood form an open set. At a limit point in \(U\), continuity
of every derivative makes its Taylor coefficients zero; its convergent
Taylor series makes the function zero on a neighborhood. Thus that set
is also relatively closed, and connectedness completes the argument.
\end{proof}

\begin{proof}[Proof of Theorem~\ref{thm:main}]
Choose injective linear maps \(A_i:\R^{m_i}\to\R^d\), with
\(m_i=\dim D_i\ge2\), such that \(D_i=A_i(B^{m_i})\), and put
\(Q_i=A_iA_i^T\). Such maps exist by restricting a singular-value
decomposition to its nonzero singular directions.
For a normal \(u\in\R^d\), the support value of \(D_i\) is
\[
 h_i(u)=\|A_i^Tu\|=\sqrt{u^TQ_i u}.
\]
When \(A_i^Tu\ne0\), equality in Cauchy--Schwarz gives the unique
maximizing point
\begin{equation}\label{eq:point}
 p_i(u)=\frac{A_iA_i^Tu}{\|A_i^Tu\|}
       =\frac{Q_i u}{\sqrt{u^TQ_i u}}.
\end{equation}
When \(A_i^Tu=0\), the entire positive-dimensional disc maximizes.

The face of a Minkowski sum exposed by \(u\) is the sum of the
individual faces: support values add, and equality in their sum is
equivalent to equality term by term. A sum of nonempty faces is a
singleton exactly when each face is a singleton. Indeed, fixing a
point in every other face translates any nonsingleton face injectively
into the sum. Therefore, with
\begin{equation}\label{eq:map}
 U=\R^d\setminus\bigcup_i\ker A_i^T,
 \qquad F(u)=\sum_i\frac{Q_i u}{\sqrt{u^TQ_i u}},
\end{equation}
we have the exact set equality \(\Exp(D)=F(U)\). The origin is excluded
from \(U\), so every normal used is nonzero.

Every excluded kernel has codimension \(m_i\ge2\).
Lemma~\ref{lem:connected} makes \(U\) connected and open. On \(U\),
every radicand is \(\|A_i^Tu\|^2>0\). The positive real inverse square
root is real analytic on \((0,\infty)\), so \(F\) is real analytic on
all of \(U\). Lemma~\ref{lem:analytic} now proves the assertion.
\end{proof}

This argument allows repeated discs, coincident spans, dependent
quadratic radicals, noninjective support parametrizations, and
lower-dimensional sums. It neither chooses a global complex
square-root branch nor assumes a rank condition on the differential
of \(F\).

\section{The generic purely nonlinear part}\label{sec:generic}

\begin{theorem}\label{thm:generic}
Suppose \(D\) is full-dimensional, every \(D_i\) has dimension at
least two, and its summand spans satisfy
\begin{equation}\label{eq:GP}
 \dim\Bigl(\sum_{i\in J}L_i\Bigr)
 =\min\Bigl(d,\sum_{i\in J}\dim L_i\Bigr)
 \quad\text{for every }J\subseteq\{1,\ldots,N\}.
 \tag{GP}
\end{equation}
Then \(S(D)=E(D)\); in particular \(S(D)\) is irreducible over \(\C\).
\end{theorem}
\begin{proof}
Every exposed point is a sum of the points \eqref{eq:point}, each on
its relative summand boundary, so \(E(D)\subseteq S(D)\).
Conversely, take \(x=\sum_i x_i\in D^\partial\cap\partial D\), where
\(x_i\in\partial_{L_i}D_i\). A supporting hyperplane at \(x\) has a
nonzero normal \(u\). Since support values add and each support deficit
is nonnegative, every \(x_i\) maximizes \(u\) on its summand.
Let \(J=\{j:A_j^Tu=0\}\). For \(i\notin J\), uniqueness gives
\(x_i=p_i(u)\).

All \(L_j\), \(j\in J\), lie in \(u^\perp\). Condition (GP) forces
their total dimension to be at most \(d-1\): otherwise their sum would
have dimension \(d\). It then gives equality between the dimension of
their sum and the sum of their dimensions. Thus the concatenated
matrix \(A_J=[A_j]_{j\in J}\) is injective and \(A_J^T\) is surjective.
Write \(x_j=A_jv_j\) with \(\|v_j\|=1\); the relative boundary
condition and injectivity justify this representation. Choose
\(w\in\R^d\) such that \(A_j^Tw=v_j\) for every \(j\in J\).

If \(J\) is empty, \(x\) is already exposed. Otherwise, for
\(\varepsilon>0\), set \(u_\varepsilon=u+\varepsilon w\). Then
\[
 A_j^Tu_\varepsilon=\varepsilon v_j,
 \qquad p_j(u_\varepsilon)=A_jv_j=x_j\quad(j\in J).
\]
For every remaining index, \(A_i^Tu_\varepsilon\) stays nonzero when
\(\varepsilon\) is sufficiently small, and
\(p_i(u_\varepsilon)\to p_i(u)=x_i\). Finiteness gives one such bound
for all indices. Consequently \(u_\varepsilon\in U\) and
\(F(u_\varepsilon)\to x\). Every point generating \(S(D)\) is a
Euclidean limit of exposed points. Since the complex algebraic set
\(E(D)\) is Euclidean closed, \(S(D)\subseteq E(D)\).
Theorem~\ref{thm:main} completes the proof.
\end{proof}

For prescribed sizes \(m_i\le d\), (GP) is a nonempty real
Zariski-open condition on the full-rank matrices \(A_i\). Each failure
of maximal concatenated rank is defined by minors. A simultaneous real
witness is obtained by partitioning distinct Vandermonde columns
\((1,t,\ldots,t^{d-1})^T\) into blocks of sizes \(m_i\). Every set of
at most \(d\) columns is independent, and every larger set has rank
\(d\). When \(\sum_i m_i\ge d\), this also ensures full dimensionality.
This matches the bases-parametrized genericity convention in
\cite[Section~2]{GM2022}. As in that paper, a smaller ambient span may
be replaced by the span of the discotope. Thus
Theorem~\ref{thm:generic} proves its Conjecture~8.2 for every feasible
generic type with no one-dimensional summands.

\section{Why the hypotheses and the distinction matter}\label{sec:limits}

One-dimensional summands cannot be admitted indiscriminately.
The interval \([-1,1]\) has exposed-point closure \(\{-1,1\}\), which
is reducible. The planar sum of the unit disc and \([-a,a]e_1\),
\(a>0\), has exposed points on two open semicircles. Its complex
Zariski closure is the union of the two distinct conics
\((X-a)^2+Y^2=1\) and \((X+a)^2+Y^2=1\).

Nor does \(S=E\) hold in arbitrary nongeneric position. In \(\R^3\),
take two identical unit discs in the \(xy\)-plane and one in the
\(xz\)-plane. The point \(e_3=e_1+(-e_1)+e_3\) belongs to
\(D^\partial\cap\partial D\), since the normal \(e_3\) exposes the
face \(2B^2_{xy}+e_3\). But \(E(D)\) lies in the zero set of
\[
 P(X,Y,Z)=(X^2+Y^2+Z^2-5)^2-4(4-Y^2)(1-Z^2).
\]
Indeed, for regular normals, write \(r^2=u_1^2+u_2^2\),
\(s^2=u_1^2+u_3^2\), \(a=2u_1/r\), and \(b=u_1/s\).
Then \(X=a+b\), \(a^2+Y^2=4\), and \(b^2+Z^2=1\), giving
\(X^2+Y^2+Z^2-5=2ab\) and \(P=0\). As \(P(0,0,1)=16\), we obtain
\(e_3\in S(D)\setminus E(D)\). This is algebraic separation, not
merely a failure of exposure by a selected normal.

\section*{Verification and use of AI}

The accompanying verification package contains the manuscript source,
exact symbolic checks with a pinned environment, source provenance,
and independent adversarial reviews. The finite checks verify the
separating polynomial, concrete general-position perturbations
(including non-orthonormal ellipses), and boundary examples. They
corroborate the formulas and do not replace the proofs of the universal
theorems above.

Generative AI tools were used extensively for proof development,
literature triage, manuscript preparation, adversarial mathematical
review, and symbolic-check design. AI review is not human peer review
or a formal proof certificate. This preprint is unrefereed. A dated
priority audit is included in the package; historical first priority
is not asserted.

\begin{thebibliography}{9}
\bibitem{GM2022}
F.~Gesmundo and C.~Meroni, \emph{The Geometry of Discotopes},
Le Matematiche \textbf{77}(1) (2022), 143--171.
\href{https://doi.org/10.4418/2022.77.1.8}{doi:10.4418/2022.77.1.8}.
\href{https://arxiv.org/abs/2111.01241}{arXiv:2111.01241}.
\bibitem{KK2017}
W.~Kucharz and K.~Kurdyka, \emph{Rationality of semialgebraic functions},
in \emph{Analytic and Algebraic Geometry 2},
Łódź University Press, Łódź, 2017, 85--96.
\href{https://doi.org/10.18778/8088-922-4.14}{doi:10.18778/8088-922-4.14}.
\bibitem{Meroni2023}
C.~Meroni, \emph{Two convex conjectures for different flavours},
in \emph{New Directions in Real Algebraic Geometry},
Oberwolfach Reports \textbf{20} (2023), Report 15/2023, 829--832;
Conjecture~1, p.~830.
\href{https://doi.org/10.4171/OWR/2023/15}{doi:10.4171/OWR/2023/15}.
\end{thebibliography}
\end{document}
